\chapter{Requirement Engineering}

\section{Overview}

This chapter describes what the system needs to do and how it should behave. Requirements
are divided into functional requirements, which describe what the system does, and
non-functional requirements, which describe how well it does it. We also describe the data
requirements and the specific rules for ambiguity detection.

\section{Stakeholder Identification}

The people who would use or be affected by this system include the following.

\textbf{Researchers and data analysts} who want to understand emotional trends in large amounts
of social media data. They need a system that is reliable and gives useful outputs, including
when the system is unsure.

\textbf{Developers} who want to integrate emotion detection into a larger application, such as a
content moderation system or a mental health monitoring tool. They need a system that is
well-documented and produces structured output.

\textbf{End users} who interact with the system directly by typing a sentence and wanting to know
its emotional content. They need clear, understandable outputs.

\section{Functional Requirements}

The following functional requirements define what the system must do.

\begin{description}[leftmargin=2.5cm, style=nextline]
    \item[FR1.] The system must accept any piece of English text as input through a command-line interface.
    \item[FR2.] The system must classify the input into one of four categories: Positive, Negative, Neutral, or Ambiguous.
    \item[FR3.] The system must convert input text into dense semantic vectors using a pre-trained sentence embedding model.
    \item[FR4.] The system must extract lexical features from input text using TF-IDF with a vocabulary of at least 3000 terms.
    \item[FR5.] The system must find the nearest neighbours of the input in the training embedding space using cosine distance.
    \item[FR6.] The system must compute Vietoris-Rips persistent homology on the neighbourhood of the input, capturing H0 and H1 features.
    \item[FR7.] The system must extract a Persistence Landscape representation from the persistence diagram.
    \item[FR8.] The system must combine semantic, lexical, and topological features into a single feature vector for classification.
    \item[FR9.] The system must use XGBoost classifiers trained in a one-vs-rest fashion for the three emotion classes.
    \item[FR10.] The system must compute the confidence gap between the top two class probabilities.
    \item[FR11.] The system must apply the ambiguity decision rules to determine the final output.
    \item[FR12.] The system must display debug information with each prediction, including the probabilities, confidence gap, loop count, and the trigger reason if the output is Ambiguous.
    \item[FR13.] The system must load a pre-trained model from disk at startup without requiring retraining.
    \item[FR14.] The system must run in an interactive loop, allowing the user to test multiple inputs in a single session.
\end{description}

\section{Non-Functional Requirements}

The following non-functional requirements describe the quality of the system.

\begin{description}[leftmargin=2.5cm, style=nextline]
    \item[NFR1.] \textbf{Performance:} The system must return a prediction in under 5 seconds per input on a standard laptop CPU.
    \item[NFR2.] \textbf{Portability:} The system must run on any machine with Python 3.10 or later installed. It must not require a GPU for inference.
    \item[NFR3.] \textbf{Reliability:} The system must handle unusual inputs, very short inputs, and empty inputs without crashing.
    \item[NFR4.] \textbf{Explainability:} When a text is flagged as Ambiguous, the system must tell the user which rule triggered the decision.
    \item[NFR5.] \textbf{Reproducibility:} Given the same input, the system must always produce the same output.
    \item[NFR6.] \textbf{Scalability:} The model must be trained once and saved. It must not require retraining to handle new inputs.
\end{description}

\section{Data Requirements}

\begin{description}[leftmargin=2.5cm, style=nextline]
    \item[DR1.] The system must use the GoEmotions dataset (simplified version) for training and evaluation.
    \item[DR2.] The 28 GoEmotions emotion categories must be mapped to three classes: Positive, Negative, and Neutral.
    \item[DR3.] The full training split of approximately 43,000 samples must be used.
    \item[DR4.] The full test split of approximately 5,400 samples must be used for evaluation.
    \item[DR5.] The trained model, the TF-IDF vectorizer, the nearest-neighbour index, and the training embeddings must all be saved to disk after training so that inference does not require retraining.
\end{description}

\section{Class Imbalance Handling Requirements}

The GoEmotions dataset is not evenly distributed across the three classes. After mapping, the
Neutral class has significantly more samples than Positive or Negative. The system must handle
this imbalance. The following requirements apply.

\begin{description}[leftmargin=2.5cm, style=nextline]
    \item[IR1.] The system must not use SMOTE or any other synthetic data generation technique.
    The reason is that SMOTE creates interpolated points in the embedding space. Since the
    topological analysis depends on the real geometric structure of the data, synthetic points
    would create artificial topology that does not correspond to real language. This would
    make the ambiguity detection unreliable.
    \item[IR2.] The system must handle imbalance through threshold tuning. Instead of using a
    fixed 0.5 decision threshold for each classifier, the threshold for each class is adjusted
    based on the class distribution. This allows the system to be more or less sensitive to
    each class without distorting the underlying data.
\end{description}

\section{Ambiguity Detection Rules}

The ambiguity detection mechanism is a central requirement of this project. The following rules
define when a text is classified as Ambiguous.

\begin{description}[leftmargin=2.5cm, style=nextline]
    \item[AR1.] If the number of H1 loops in the local persistence diagram is 14 or more, the
    text is Ambiguous regardless of the classifier probabilities. This is called the
    \textit{extreme geometry condition}.
    \item[AR2.] If the number of H1 loops is 10 or more, and the confidence gap between the
    top two predicted classes is less than 0.30, the text is Ambiguous. This is called the
    \textit{moderate geometry combined with low confidence condition}.
    \item[AR3.] If the confidence gap between the top two predicted classes is less than 0.10,
    regardless of the loop count, the text is Ambiguous. This is called the
    \textit{very low confidence condition}.
    \item[AR4.] If none of the above conditions are met, the text receives the label of the class
    with the highest probability.
\end{description}

The thresholds (14, 10, and 0.30) were selected based on experimentation. The loop count
thresholds were chosen by looking at the distribution of loop counts across samples and
identifying values that separated clearly ambiguous inputs from clearly unambiguous ones. The
confidence gap threshold of 0.30 was chosen because below this value, the classifier output is
not stable enough to be trusted.

\section{Technology Stack}

\textbf{Table~\ref{tab:techstack}} lists the key technologies used to build the system.

\begin{table}[H]
\centering
\caption{Technology Stack}
\label{tab:techstack}
\begin{tabular}{lll}
\toprule
\textbf{Category} & \textbf{Technology} & \textbf{Purpose} \\
\midrule
Programming Language & Python 3.10+        & Core implementation \\
Machine Learning     & XGBoost             & One-vs-Rest classification \\
NLP Embeddings       & SentenceTransformers & Semantic embeddings (all-mpnet-base-v2) \\
Feature Extraction   & scikit-learn        & TF-IDF vectorization, KNN search \\
Topological Analysis & Ripser              & Vietoris-Rips persistent homology \\
Topology Processing  & Persim              & Persistence Landscape computation \\
Data Handling        & NumPy, Pandas       & Array operations, data manipulation \\
Dataset              & GoEmotions          & Training and evaluation data \\
Model Persistence    & Joblib              & Saving and loading trained models \\
Development IDE      & VS Code             & Code editing and debugging \\
\bottomrule
\end{tabular}
\end{table}

The system is built entirely in Python and does not require a GPU for inference. All dependencies
are installable via \texttt{pip} and are compatible with standard laptops and desktops running
Windows, macOS, or Linux.

\section{Requirements Summary}

\textbf{Table~\ref{tab:req_summary}} summarizes the key requirements of the system.

\begin{table}[H]
\centering
\caption{Summary of Key System Requirements}
\label{tab:req_summary}
\begin{tabular}{llp{7cm}l}
\toprule
\textbf{ID} & \textbf{Type} & \textbf{Requirement} & \textbf{Priority} \\
\midrule
FR1  & Functional     & Accept English text input                              & High   \\
FR2  & Functional     & Output Positive / Negative / Neutral / Ambiguous       & High   \\
FR3  & Functional     & Compute local topology at inference                    & High   \\
FR4  & Functional     & Apply ambiguity decision rules                         & High   \\
FR5  & Functional     & Display debug information per query                    & Medium \\
FR13 & Functional     & Load pre-trained model from disk                       & High   \\
NFR1 & Non-Functional & Respond in under 5 seconds                             & High   \\
NFR2 & Non-Functional & Run without GPU for inference                          & High   \\
NFR4 & Non-Functional & Explain why text is flagged as Ambiguous               & Medium \\
IR1  & Imbalance      & Do not use SMOTE or synthetic data generation          & High   \\
IR2  & Imbalance      & Use threshold tuning for imbalance handling            & High   \\
AR1  & Ambiguity      & Flag on extreme loop count ($\geq 14$)                 & High   \\
AR2  & Ambiguity      & Flag on moderate topology + low confidence             & High   \\
AR3  & Ambiguity      & Flag on very low confidence gap ($< 0.10$)             & High   \\
\bottomrule
\end{tabular}
\end{table}

As shown in \textbf{Table~\ref{tab:req_summary}}, high priority is given to the core classification and
ambiguity detection mechanisms, while user interface and debug features are treated as medium
priority.
